\chapter{Background and Related Work}\label{ch:background}

This chapter establishes the foundations on which the thesis builds and situates
the work against the prior art. It draws on four largely separate lines of
research---deductive program verification, specification inference,
large-language-model (LLM) software engineering, and the empirical study of
software evolution---and argues that none, alone, spans the problem addressed here.
The survey is critical rather than encyclopaedic, naming what each tradition
achieves and where it breaks down, so that the gap emerges by subtraction and
\cref{sec:bg-synthesis} can state the thesis's position within it.

\section{Design by contract and its axiomatic antecedents}\label{sec:bg-dbc}

A \emph{specification} is a statement of intended behaviour independent of any
implementation. The formal-methods tradition renders it as a predicate, or a pair
of predicates, over program states, and every contract here descends from it. The
axiomatic foundation is Hoare logic, in which the triple $\{P\}\,S\,\{Q\}$ asserts
that if precondition $P$ holds before command $S$ executes and $S$ terminates, then
postcondition $Q$ holds afterward \cite{hoare1969axiomatic}. As Hoare defined it
the triple expresses \emph{partial} correctness---it promises nothing about
termination, only that \emph{if} $S$ halts, $Q$ holds---a distinction that recurs
throughout, since the extended static checking used here is a partial-correctness
regime. Dijkstra recast this as a calculation: the weakest precondition
$\mathrm{wp}(S,Q)$ is the weakest predicate on the initial state guaranteeing that
$S$ terminates and establishes $Q$ \cite{dijkstra1976discipline,dijkstra1975guarded}.
The calculus matters structurally as well as historically: because $\mathrm{wp}$
is a computable function of the postcondition and the statement, a change to the
postcondition a callee guarantees can be pushed back mechanically through its call
sites---the technical heart of the compositional change-propagation analysis of
\cref{ch:compositional}.

The refinement calculus extended this into a development method that transforms an
abstract specification into code by provably correct steps
\cite{back1998refinement,morgan1994programming}---the ancestor of program synthesis
and, through it, of the contract-to-code generation LLM agents now perform. The
decisive difference is that refinement \emph{trusts the derivation}, each step
justified by a proof rule, whereas an LLM agent offers no such justification. The
thesis therefore inverts the relationship, letting a fallible generator propose an
implementation freely and recovering soundness by \emph{verifying} the result
against the contract rather than trusting the process that produced it.

Design by contract, introduced by Meyer for Eiffel, operationalised these
semantics for object-oriented software \cite{meyer1992dbc,meyer1997oosc}. A method
is governed by a precondition the caller must establish, a postcondition the method
guarantees in return, and class invariants every public method must preserve. The
conceptual contribution is the framing of behaviour as mutual \emph{obligations}
between caller and callee: the caller must satisfy the precondition and benefits
from the postcondition; the callee may assume the precondition and must deliver the
postcondition. This duality is foundational, because it is exactly what makes
contracts \emph{compose}---a change to a callee's contract recomputes every
caller's obligations, the property the compatibility analysis exploits. It carries
a further consequence: if the contract is the authoritative statement of intent and
the code merely its fulfilment, then when code can be regenerated cheaply the
contract, not the code, is the scarce and durable asset. Eiffel's limitation is its
medium---assertions are evaluated at runtime against whatever inputs occur, so the
guarantee is sampled rather than universal.

The behavioural interface specification tradition closes that gap. The Java
Modeling Language (JML) is the contract notation adopted throughout
\cite{leavens2006jml}. Its annotations live in specially marked comments,
transparent to the ordinary compiler, and use a superset of Java's expression
syntax: a \texttt{requires} clause states a precondition, \texttt{ensures} a
postcondition, \texttt{invariant} a class invariant, and \texttt{assignable} (the
frame clause) bounds the locations a method may modify---indispensable for sound
modular reasoning because it tells a verifier what a call cannot change.
Postconditions may refer to the pre-state through \verb|\old(e)|, to the return
value through \verb|\result|, and may quantify with \verb|\forall| and
\verb|\exists|. Two properties are load-bearing. A JML contract is
\emph{machine-checkable}: unlike a natural-language comment it has a formal
semantics against which an implementation can be judged mechanically, which
qualifies it as an oracle rather than documentation. And a contract is
\emph{partial}: it constrains behaviour only so far as its clauses reach, which lets
inference build a contract incrementally but makes a vacuous contract a first-class
hazard---trivially satisfied, and so excusing incorrect code---that the thesis
actively detects. The closest contemporary relative, Spec\#, embeds contracts in
C\# with a verifying compiler \cite{barnett2005specsharp}; JML's tooling maturity
for Java makes it the practical substrate here.

\section{Verification: OpenJML, extended static checking, and SMT}\label{sec:bg-verification}

A machine-checkable contract is useful only if an implementation can be checked
against it across all inputs. Deductive verification reduces ``does this code
satisfy this contract?'' to a collection of \emph{verification conditions} whose
validity entails correctness, produced by a weakest-precondition computation over
the method body \cite{dijkstra1976discipline,manna1980deductive}. Loops are where
the backward walk breaks down, since an unknown number of iterations cannot be
enumerated; the remedy is a \emph{loop invariant} the verifier assumes in place of
unrolling. An invariant must be both \emph{inductive}---true on entry and preserved
by the body---and \emph{strong enough} that, with the negated guard at exit, it
entails the required obligation. Supplying invariants that meet both conditions is
the largest source of the annotation burden that has historically deterred
adoption, and a principal target of the inference surveyed in
\cref{sec:bg-inference}.

Extended static checking (ESC), pioneered by ESC/Java \cite{flanagan2002escjava},
is the pragmatic engineering of this reduction: it sacrifices completeness for
automation and modularity, analysing each method in isolation while \emph{assuming}
the contracts of its callees rather than inlining their bodies. Modularity is what
makes ESC scale and makes contract composition meaningful---a proof depends only on
callee contracts, not implementations---but it bounds the guarantee: an ESC proof
is only as sound as the contracts it assumes. OpenJML is the modern realisation of
ESC for Java built on JML \cite{cok2011openjml,cok2014openjml}; this thesis adopts
it as a trusted oracle rather than proposing a new prover. Beneath it is the Z3 SMT
solver \cite{demoura2008z3}, deciding satisfiability of first-order formulas over
the background theories---linear arithmetic, arrays, bit-vectors, uninterpreted
functions---that Java constructs generate. OpenJML poses each condition in negated
form: \texttt{unsat} certifies validity, while \texttt{sat} yields a
\emph{counterexample}, a concrete witness of violation that the integration chapter
exploits as a signal far more informative than the bare fact that a test failed.
Under modular checking, however, such a witness can be \emph{spurious}, reflecting
a too-weak invariant rather than a reachable defect, and the thesis does not
overclaim what it establishes.

Precision about what a proof guarantees is essential. A successful ESC run proves
the implementation satisfies its contract for \emph{all} admitted inputs---a
guarantee no finite test suite provides---but does \emph{not} prove the contract
captures what the developer wanted; that is the province of inference, not the
verifier. The decidability limits must also be stated honestly. Verification is
undecidable in general, and although Z3's quantifier-free fragments are decidable,
the quantified formulas JML routinely generates are handled by heuristic
instantiation and may return \texttt{unknown} or time out---in practice on
nonlinear arithmetic, rich array reasoning, and deeply nested quantifiers. The
corpora here exhibit exactly these modes, reported as solver \emph{incompleteness
and timeout}, not unsoundness in the inferred specifications. These limits are
intrinsic to the condition-to-SMT architecture rather than artefacts of one tool,
as the wider verifier ecosystem---Dafny \cite{leino2010dafny}, KeY
\cite{ahrendt2016key}, Frama-C \cite{kirchner2015framac}, Verus
\cite{lattuada2023verus}, and the cvc5 backend \cite{barbosa2022cvc5}---confirms. A
parallel synthesis tradition, counterexample-guided inductive synthesis and
sketching \cite{solarlezama2006sketching,solarlezama2008sketching} standardised as
syntax-guided synthesis \cite{alur2013sygus}, closes a generate--check--refine loop
against a machine-checkable oracle---precisely the shape \cref{ch:integration} gives
its synthesis--verification loop, with an LLM in place of enumerative search.

\section{Specification inference and the oracle problem}\label{sec:bg-inference}

The chief obstacle to deploying any contract-based machinery is the cost of
\emph{authoring} contracts. Specification inference attacks this bootstrap cost and
is the technical core of the thesis's first claim, that specifications can be made
cheap. The dynamic tradition observes a running program and reports the properties
that held. Its canonical instrument, Daikon, records variable values at chosen
points and reports the templated properties that survived every observed trace
\cite{ernst2007daikon}, establishing that a specification can be \emph{recovered}
from behaviour rather than authored from a blank page; specification mining
similarly learns finite-state API protocols from traces \cite{ammons2002mining}.
The defining limitation is soundness: a mined property is only \emph{likely} true,
summarising sampled executions rather than all of them, so integrating it with a
static checker is itself a research thread \cite{nimmer2001static}. Such an
invariant describes what the code \emph{did}, not what it \emph{should} do, and so
cannot distinguish a faithful contract from an encoded bug.

The static tradition buys soundness by over-approximation. Abstract interpretation
computes fixpoints over abstract domains so that every reported property provably
holds \cite{cousot1977abstract}; Houdini takes the complementary guess-and-check
route, postulating candidate annotations and using an ESC verifier to refute the
unprovable ones \cite{flanagan2001houdini}---the historical antecedent of the
verifier-in-the-loop discipline that recurs throughout this thesis. The price of
soundness is real: rich domains do not always converge, and the decision
procedure's limits bound what can be discharged. A learning-based strand poses
invariant synthesis as learning from examples and verifier-supplied counterexamples
\cite{garg2014ice,si2018learning}, and since 2023 the focus has shifted to LLMs
prompted to emit candidate invariants a symbolic tool then checks, as in Loopy
\cite{kamath2023loopy}, or extended to full JML specifications with mutation-based
fallback, as in SpecGen \cite{ma2025specgen}; a separate sub-thread translates
natural language directly into formal artefacts
\cite{cosler2023nl2spec,endres2024nl2postcond,lahiri2025nl2contract}. The common
commitment the thesis endorses is that the generative model is never trusted alone
but paired with a sound checker admitting only what it can prove.

The instrument carried through this thesis, the JML-Inferrer, occupies the static
end of this space. It is a Java~21 tool over the JavaParser abstract syntax tree
that emits JML---preconditions, postconditions, loop and class invariants, and
\texttt{assignable} clauses---by heuristic pattern matching over syntax rather than
complete program analysis: a null-check induces a non-nullness precondition,
recognised loop shapes induce matching invariants and result postconditions, and
numeric guards induce range and overflow preconditions. The heuristic basis is a
deliberate trade, favouring recall of plausible, valid JML over the soundness
guarantees of abstract interpretation, on the premise that the checker will reject
any clause that does not hold---so an unsound guess costs a failed proof, not a
false guarantee. The inferrer is thus static like Houdini but generative rather than
refutational, over a rich JML vocabulary. Throughout, the thesis distinguishes a
\emph{verified} clause, one OpenJML has discharged, from a merely \emph{well-formed}
clause, one that is syntactically valid and scope-correct but unchecked, and never
lets the distinction blur.

The dual of inference is the oracle problem. A test exercises a program; an
\emph{oracle} decides whether the observed behaviour is correct.
\citet{barr2015oracle} establish that while input generation is now largely
automatable, the oracle remains the binding constraint, distinguishing
\emph{specified} oracles, derived from a formal description, from \emph{derived} and
\emph{implicit} ones. Search-based generators sharpen the gap: lacking a
specification, EvoSuite emits regression assertions that merely capture current
behaviour \cite{fraser2011evosuite} and Randoop flags only crashes and contract
violations \cite{pacheco2007randoop}. LLMs have accelerated this without resolving
it: TestPilot repairs tests by re-prompting on failures \cite{schafer2024testpilot},
but recent analysis documents the hazard the thesis takes as its motivating
defect---when the same model both implements the code and writes the asserting
oracle, the oracle encodes the model's own mistaken expectations, yielding tests
that pass by construction \cite{liu2026oracle}. This thesis instead
elevates an inferred, verified JML contract to a specified oracle discharged by
OpenJML over all inputs, separating the authority that writes the contract from the
agent that satisfies it.

\section{LLM code generation and spec-driven lifecycles}\label{sec:bg-llm}

The generators the verification machinery is meant to gate are LLMs and the agents
built on them. Codex demonstrated function-level generation from docstrings and
introduced the \emph{pass@k} metric and HumanEval \cite{chen2021codex}, with MBPP a
complementary benchmark \cite{austin2021mbpp} and AlphaCode reaching median human
performance in contests \cite{li2022alphacode}. These systems inverted the
economics of classical synthesis---natural-language intent is cheap where formal
specifications were not---but did so by abandoning the closed loop: correctness is
\emph{estimated} against a finite, model-adjacent test suite rather than proved.
That estimate is fragile in two ways. Test suites are frequently too weak to expose
plausible-but-wrong solutions, as EvalPlus showed by lowering reported pass rates
through added tests \cite{liu2023evalplus}; and generation is prone to
\emph{hallucination}, the confident fabrication of non-existent APIs and subtly
inconsistent logic \cite{zhang2025hallucination}. The deeper problem is that
generator and informal oracle share a model and training distribution, so
verification is effectively self-grading, even as LLMs reshape the lifecycle
wholesale \cite{ozkaya2023llmse}.

The frontier has moved from isolated functions to \emph{agents} over whole
repositories. SWE-bench reframed evaluation around resolving real GitHub issues
with project test suites as the oracle \cite{jimenez2024swebench}, and SWE-agent
showed a purpose-built interface markedly improves navigation and editing
\cite{yang2024sweagent}. The dominant oracle across this work remains test
execution, inheriting the incompleteness of testing and, when tests are
model-authored, the self-grading circularity at repository scale. Most directly
related is the emerging family of \emph{spec-driven}, AI-native lifecycles, which
reorganise the process so a specification, not code, is the primary human-authored
artefact---consonant with \cref{sec:bg-dbc}, that when an agent regenerates the
implementation cheaply the durable asset is the specification. GitHub's Spec Kit
formalises a specify--plan--tasks--implement workflow anchored by a project
``constitution'' \cite{githubspeckit}; Amazon's AI-Driven Development Lifecycle
restructures the lifecycle into inception, construction, and operations phases with
quality gates and a brownfield reverse-engineering step \cite{awsaidlc}; and the
Kiro IDE places a requirements-to-design-to-tasks workflow at the editor's centre
\cite{awskiro}. These efforts correctly identify the economic inversion but share
three structural defects the integration chapter remedies. The specification is
\emph{natural language}, so the specification-to-code step is a compilation with no
checker relating generated code to stated intent. The loop is \emph{open}:
generation is followed by review-by-eye, not verification. And where verification
appears at all it is \emph{self-grading} test execution. A natural-language
artefact is, besides, too imprecise to support automated reasoning over
change---these lifecycles cannot compute which dependents a modification breaks
without re-running anything.

\section{Version compatibility as an implicit-contract problem}\label{sec:bg-evolution}

Modern software is assembled from libraries, and the contract between a library and
its clients is communicated by a version number under semantic versioning: a major
increment signals an incompatible change, a minor one backward-compatible
additions, a patch a backward-compatible fix \cite{preston2013semver}. The promise
is honoured unevenly. \citet{raemaekers2017semver} found roughly a third of Maven
Central releases introduce a breaking change, frequently without the major
increment; \citet{ochoa2022breaking} replicated this over a far larger corpus,
reporting higher compliance but noting most breaking changes touch code no client
uses; and \citet{decan2019empirical} characterise backward-incompatible updates as
a structural hazard across packaging ecosystems. The version string is thus a weak
proxy for what developers actually care about---whether \emph{their} call sites
still behave correctly.

Signature-diffing tools such as APIDiff compute deltas of type-, method-, and
field-level changes \cite{brito2018apidiff}, and large studies catalogue such
changes and their causes \cite{xavier2017historical,brito2018why}, but share a
ceiling: they detect changes to the \emph{shape} of an interface, not its
\emph{behaviour}. A method whose signature is untouched but whose contract has
silently shifted---an off-by-one boundary, a newly thrown exception---is invisible.
This motivates work on \emph{behavioural} incompatibility, the closest prior art to
the compatibility chapters: \citet{mostafa2017behavioral} showed such
signature-preserving changes are common and disproportionately damaging precisely
because no diffing tool flags them, \citet{jezek2017magic} that detection rates
vary widely across tools, and \citet{jayasuriya2024understanding} quantified breaks
detectable only when client tests run. The recurring move is to recover behaviour
\emph{dynamically}---through client tests or mined invariants---which observes
sampled inputs and so confirms an incompatibility but never its absence; a recent
LLM direction \emph{repairs} breaks once detected \cite{islam2025byam} but
presupposes detection and guarantees nothing about the repair. The decisive gap is
that compatibility is treated as an \emph{ex post}, sampled, largely syntactic
property; none of these strands answers the question posed at upgrade time---which
of my call sites does this change break, and why---with a proof obligation rather
than a failed test.

\section{Positioning: the gap this thesis addresses}\label{sec:bg-synthesis}

Each line of work optimises one edge of the problem in isolation. The verification
community possesses a machine-checkable contract layer \cite{cok2014openjml,leino2010dafny}
but presupposes a human authors the contract. The inference tradition recovers
contracts cheaply \cite{ernst2007daikon,flanagan2001houdini} but treats the
inferred property as a terminal deliverable that describes what the code does, not
what it should do. The LLM and agentic literature solved the authoring-cost problem
that crippled classical synthesis \cite{chen2021codex,yang2024sweagent} but in
doing so discarded the machine-checkable contract and the closed loop that made
synthesis trustworthy, its spec-driven lifecycles institutionalising three
defects---natural-language specs, an open loop, and self-grading verification
\cite{githubspeckit,awsaidlc,liu2026oracle}. And the behavioural-compatibility
literature measures incompatibilities after the fact
\cite{ochoa2022breaking,islam2025byam} but does not derive a change's blast radius
from a verified interface without running a client.

No prior line combines these capabilities, and it is the conjunction, not any
single ingredient, that is novel. This thesis differs by being static,
syntax-driven, and verification-disciplined: contracts are inferred from source
without test inputs or traces, deterministically from the syntax, and every clause
is held to OpenJML discharge or reported as a solver limit rather than shipped as an
approximation. It is \emph{compositional}: because contracts compose along the
caller/callee obligation structure of design by contract, a callee-spec change
recomputes each caller's proof obligations and flags behaviourally broken call
sites without execution, making an internal refactor and a third-party version bump
the same static blast-radius operation (\cref{ch:compositional}). And it is
\emph{integration-oriented}: the inferred, verified contract is interposed as the
pivot of an AI-native lifecycle---the input an agent builds against and the oracle
its output is verified against over all inputs---authored by an authority that never
grades its own output (\cref{ch:integration}). \Cref{ch:inferrer} establishes that
such cheap, partial inference is nonetheless \emph{useful}, measurably improving
LLM-generated tests. The remaining chapters develop these capabilities in turn,
reusing the inference and verification machinery established here as trusted
components, and returning in \cref{ch:discussion,ch:conclusion} to the limits that
bound them---above all the boundary between verified and merely well-formed
specifications.
